%%%%%%%%%%%%%%%%%%%%%%%%%%%%%%%%%%%%%%%%%%%%%%%%%%
\needspace{12\baselineskip}
\section{Boolesche Algebra}
%%%%%%%%%%%%%%%%%%%%%%%%%%%%%%%%%%%%%%%%%%%%%%%%%%
\worum
In der Booleschen Algebra (nach George Boole) haben Variablen nur zwei Werte: \textbf{0 (falsch)} und \textbf{1 (wahr)}. Mit wenigen Verknüpfungen lassen sich daraus beliebig komplizierte Bedingungen bauen. Jede Bedingung in einem Programm (\code{if}, \code{while}) ist ein boolescher Ausdruck. Um zu prüfen, was ein Ausdruck liefert oder ob zwei Ausdrücke gleichwertig sind, geht man in einer \textbf{Wahrheitstafel} alle Belegungen systematisch durch. Bei $n$ Variablen sind das $2^n$ Zeilen.

\begin{merke}
\begin{center}
\begin{tabular}{cc|cccc}
$a$ & $b$ & $a \wedge b$ (AND) & $a \vee b$ (OR) & $a \oplus b$ (XOR) & $\neg a$ (NOT) \\ \hline
0 & 0 & 0 & 0 & 0 & 1 \\
0 & 1 & 0 & 1 & 1 & 1 \\
1 & 0 & 0 & 1 & 1 & 0 \\
1 & 1 & 1 & 1 & 0 & 0
\end{tabular}
\end{center}
\textbf{Wichtige Gesetze} (jeweils auch mit vertauschten Rollen von $\wedge$ und $\vee$):

\smallskip
\begin{tabular}{@{}ll@{\qquad}ll@{}}
Kommutativ: & $a \wedge b = b \wedge a$ & Neutralität: & $a \wedge 1 = a$, \ $a \vee 0 = a$ \\
Assoziativ: & $(a \wedge b) \wedge c = a \wedge (b \wedge c)$ & Extremal: & $a \wedge 0 = 0$, \ $a \vee 1 = 1$ \\
Distributiv: & $a \wedge (b \vee c) = (a \wedge b) \vee (a \wedge c)$ & Komplementär: & $a \wedge \neg a = 0$, \ $a \vee \neg a = 1$ \\
Idempotenz: & $a \wedge a = a$ & Doppelnegation: & $\neg(\neg a) = a$ \\
Absorption: & $a \vee (a \wedge b) = a$ & De Morgan: & $\neg(a \wedge b) = \neg a \vee \neg b$ \\
 & $a \wedge (a \vee b) = a$ & & $\neg(a \vee b) = \neg a \wedge \neg b$
\end{tabular}
\end{merke}

\beispiele
\beispiel{Gleichwertigkeit mit der Wahrheitstafel prüfen}
Liefern $a \vee (\neg a \wedge b)$ und $a \vee b$ für jede Belegung denselben Wert?
\begin{center}
\begin{tabular}{cc|cc|c|c}
$a$ & $b$ & $\neg a$ & $\neg a \wedge b$ & $a \vee (\neg a \wedge b)$ & $a \vee b$ \\ \hline
0 & 0 & 1 & 0 & \textbf{0} & \textbf{0} \\
0 & 1 & 1 & 1 & \textbf{1} & \textbf{1} \\
1 & 0 & 0 & 0 & \textbf{1} & \textbf{1} \\
1 & 1 & 0 & 0 & \textbf{1} & \textbf{1}
\end{tabular}
\end{center}
Die beiden letzten Spalten stimmen in allen vier Zeilen überein, die Ausdrücke sind also gleichwertig. Ein einziger Unterschied würde genügen, um das Gegenteil zu zeigen.

\beispiel{Vereinfachen mit den Gesetzen}
\begin{align*}
a \wedge \neg(\neg a \vee b)
&= a \wedge (\neg\neg a \wedge \neg b) && \text{De Morgan} \\
&= a \wedge (a \wedge \neg b) && \text{Doppelnegation} \\
&= (a \wedge a) \wedge \neg b && \text{Assoziativ} \\
&= a \wedge \neg b && \text{Idempotenz}
\end{align*}

\beispiel{Vom Text zum Ausdruck}
Ein Kunde bekommt Rabatt, wenn er Mitglied ist ($m$) oder wenn er für mehr als 100 Euro einkauft ($e$) und keine Aktionsware im Korb hat ($k$). Als Ausdruck: $r = m \vee (e \wedge \neg k)$. In Java schreibt man dafür \code{m || (e \&\& !k)}.

\aufgaben
\aufgabe{Wahrheitstafel mit drei Variablen}
Liefern die Ausdrücke $a \wedge (b \vee c)$ und $(a \wedge b) \vee c$ für jede Belegung von $a$, $b$ und $c$ denselben Wert? Überprüfen Sie dies mit einer Wahrheitstafel und begründen Sie Ihre Antwort.

\aufgabe{De Morgan}
\begin{enumerate}[label=\alph*)]
\item Weisen Sie beide Gesetze von De Morgan mit einer Wahrheitstafel nach.
\item Ein häufiger Fehler: Liefern $\neg(a \wedge b)$ und $\neg a \wedge \neg b$ dasselbe? Geben Sie eine Belegung an, die den Unterschied zeigt.
\end{enumerate}

\aufgabe{Vereinfachen}
Vereinfachen Sie mithilfe der Gesetze so weit wie möglich. Geben Sie in jedem Schritt das verwendete Gesetz an.
\begin{multicols}{2}
\begin{enumerate}[label=\alph*)]
\item $a \wedge \neg(b \vee \neg a)$
\item $\neg(a \wedge \neg(b \vee a))$
\item $(a \wedge b) \vee (a \wedge \neg b)$
\item $\neg(\neg a \vee \neg b) \vee (a \wedge c)$
\end{enumerate}
\end{multicols}

\aufgabe{Vom Text zum Ausdruck}
\begin{enumerate}[label=\alph*)]
\item Eine Heizung soll laufen, wenn es kälter als 18 Grad ist ($k$) und das Fenster nicht offen ist ($f$), oder wenn der Frostschutz aktiv ist ($s$). Geben Sie einen booleschen Ausdruck an und stellen Sie die Wahrheitstafel auf.
\item Eine Alarmanlage soll auslösen, wenn sie scharf geschaltet ist und mindestens einer von zwei Sensoren anspricht. Geben Sie einen Ausdruck an.
\end{enumerate}

\aufgabe[\S]{Mit weniger auskommen}
\begin{enumerate}[label=\alph*)]
\item Drücken Sie $a \oplus b$ (XOR) nur mit $\wedge$, $\vee$ und $\neg$ aus und prüfen Sie Ihr Ergebnis mit einer Wahrheitstafel.
\item Die Verknüpfung NAND ist definiert als $\neg(a \wedge b)$. Zeigen Sie, dass sich $\neg a$, $a \wedge b$ und $a \vee b$ allein mit NAND ausdrücken lassen.
\item Wie viele verschiedene Verknüpfungen von zwei Variablen (also verschiedene Ergebnisspalten einer Wahrheitstafel) gibt es insgesamt?
\end{enumerate}

\ueben
\aufgabe{Wahrheitstafeln}
Stellen Sie für jeden Ausdruck die Wahrheitstafel auf.
\begin{multicols}{4}
\begin{enumerate}[label=\alph*)]
\item $\neg a \vee b$
\item $a \wedge \neg b$
\item $(a \vee b) \wedge \neg(a \wedge b)$
\item $a \vee (b \wedge c)$
\end{enumerate}
\end{multicols}

\aufgabe{Auswerten}
Welchen Wert haben die Ausdrücke für $a = 1$, $b = 0$ und $c = 1$?
\begin{multicols}{4}
\begin{enumerate}[label=\alph*)]
\item $a \wedge b \vee c$
\item $\neg(a \vee b) \wedge c$
\item $a \oplus c$
\item $\neg a \vee \neg b \vee \neg c$
\end{enumerate}
\end{multicols}

\aufgabe{Gleichwertig?}
Prüfen Sie mit einer Wahrheitstafel, ob die beiden Ausdrücke gleichwertig sind.
\begin{enumerate}[label=\alph*)]
\item $a \wedge (a \vee b)$ \ und \ $a$
\item $\neg a \vee b$ \ und \ $\neg(a \wedge \neg b)$
\item $a \oplus b$ \ und \ $(a \vee b) \wedge (\neg a \vee \neg b)$
\item $(a \wedge b) \vee c$ \ und \ $(a \vee c) \wedge (b \vee c)$
\end{enumerate}

\aufgabe{Vereinfachen}
\begin{multicols}{2}
\begin{enumerate}[label=\alph*)]
\item $a \wedge (\neg a \vee b)$
\item $(a \vee b) \wedge (a \vee \neg b)$
\item $\neg(\neg a \wedge \neg b)$
\item $(a \wedge b) \vee (\neg a \wedge b) \vee (a \wedge \neg b)$
\end{enumerate}
\end{multicols}

\aufgabe{De Morgan in Java}
Formen Sie die Bedingungen mit den Gesetzen von De Morgan so um, dass kein \code{!} mehr vor einer Klammer steht.
\begin{multicols}{2}
\begin{enumerate}[label=\alph*)]
\item \code{!(x > 5 \&\& y < 3)}
\item \code{!(fertig || !gueltig)}
\end{enumerate}
\end{multicols}
